\documentclass{article}
\usepackage{kotex}

\begin{document}
\title{2026-2 `통사론(001)' 스터디 INFO}
\author{언어학과 이승헌(2025-17960)}
\date{\today}
\maketitle
\tableofcontents

\section{개요 \emph{Prologue}}

\section{계획 \emph{details}}\footnote{스터디 계획은 추후 논의를 통해 변경될 수 있음.}
\begin{itemize} [-]
\item 스터디 교재: Syntax: A Generative Introduction (4th Edition.) \textit{by Andrew Carnie}
\item 스터디 기간 및 일정: 7월 3주차(7/13)~$\sim$~8월 4주차(8/30), 추후 일정 수합 예정
\item 스터디 방식: 매주 정해진 양의 교재\footnote{매주 읽을 교재의 양은 지난 학기 통사론의 수업 진도를 참고하여 적절한 수준으로 설정할 계획임.}를 읽는 것을 기반으로. 기본적인 메커니즘 자체는 자기주도학습이나, 주 1~2회가량 모여서 함께 교재를 읽고 상호 질문 및 피드백할 것을 제안
\item 스터디 목표: 
\begin{quote}
    `교재의 내용을 빠짐없이 마스터하자!' (X) \\
    `교재를 읽으며 통사론을 가볍게나마 경험하고, 추후 학습 과정에서 교재를 \textbf{적극적으로 활용할 수 있도록} 하자!' (O)
\end{quote}
\item 스터디 구성원: 언어25 김지수, 언어25 이승헌, 언어26 류지민, 언어26 유서영
\item 스터디 장소(가안): 관정도서관 그룹스터디룸 
\end{itemize}

\section{마무리 \emph{Epilogue}}
야구도 이기고 이래저래 기분이 좋아서 \LaTeX 을 활용한 컨셉질을 좀 해보았습니다. 
하루이틀도 아니고 그러려니 해주시기 바랍니다. 
여전히 미숙한 실력이지만 그래도 이정도 문서는 만들 수 있다는 사실에 스스로 놀랍네요. 
\textbf{기기 정공 화이팅~} 
\end{document}

